% 主要研究内容与论文组织

\section{主要研究内容与创新点}

\subsection{一维极度耦合：基于紧凑坐标基变换的细长杆仿真}

% TODO: 补充紧凑坐标表示方法的贡献
本节介绍基于紧凑坐标表示的不可伸长 Cosserat 杆仿真方法。

\subsection{三维局部畸变：基于算子自适应基函数的病态网格处理}

% TODO: 补充算子自适应基函数方法的贡献
本节介绍有限元方法中病态网格单元的处理方法。

\subsection{宏观全局耦合：基于多重划分模态基函数的大规模特征值求解}

% TODO: 补充多重划分模态综合法的贡献
本节介绍基于多重划分的模态综合法求解大规模特征值问题。

\section{本文的组织结构}

本文的组织结构安排如下：
\inputbody{intro/3.1_structure}

\begin{itemize}
  \item \textbf{第一章：绪论。} 阐述物理仿真中软硬耦合系统面临的通用数值挑战，系统性地提出基于基函数与坐标变换的破局理论，并总结了本文在三个不同维度/尺度下的主要研究贡献。
  \item \textbf{第二章：基于紧凑表示的不可伸长 Cosserat 杆高效稳定仿真。} 详细论述针对一维细长杆极度耦合问题所设计的无约束坐标基函数变换方法及对应的预条件求解器设计。
  \item \textbf{第三章：有限元方法中病态网格单元的处理：数值感知细化与聚合。} 针对三维连续体局部几何畸变导致的刚性锁定，介绍算子自适应插值基底的优化与凝聚策略。
  \item \textbf{第四章：基于多重划分的模态综合法求解大规模特征值问题。} 针对宏观多组件系统的强界面耦合，推导相位互补的模态基降维技术以及免 Schur 补的高效计算框架。
  \item \textbf{第五章：总结与展望。} 对全文的基函数变换方法进行系统性回顾，总结其适用范围与优势，并对该理论在流固耦合及更广义多物理场仿真中的未来应用提出展望。
\end{itemize}
